
\section{UV Mapping tradicional}

O UV mapping estabelece uma correspondência entre a superfície de um objeto tridimensional e uma imagem bidimensional. As coordenadas \(U\) e \(V\) normalmente são normalizadas no intervalo \([0,1]\). Dessa maneira, um vértice do modelo pode carregar, além de sua posição e outros atributos, uma coordenada que indica qual região da textura deve ser associada àquele ponto.

No pipeline de renderização, as coordenadas dos vértices são fornecidas ao estágio de vértice e, durante a rasterização, são interpoladas para os fragmentos. O fragment shader pode então utilizar as coordenadas interpoladas para consultar uma textura por meio de um sampler.

Heckbert \cite{heckbert1986} destaca dois aspectos centrais do texture mapping: o mapeamento geométrico, responsável por relacionar a textura à superfície, e a filtragem, necessária para reduzir artefatos como aliasing.

\section{UV Mapping dinâmico}

O UV mapping dinâmico consiste em modificar as coordenadas de textura durante a execução. Uma transformação simples é a translação:

\[
u'=u+\Delta u,\qquad v'=v+\Delta v.
\]

Se \(\Delta u\) ou \(\Delta v\) variarem com o tempo, a textura pode apresentar movimento contínuo. Esse mecanismo pode ser empregado para representar, por exemplo, o deslocamento visual de água, energia, lava ou elementos de interface.

Também podem ser aplicadas escala e rotação. A combinação dessas operações permite controlar a velocidade, direção, repetição e orientação de uma textura sem modificar a malha original.

\section{Mapeamento procedural}

No mapeamento procedural, o resultado é produzido por uma função ou procedimento. Dong et al. \cite{dong2020} definem a geração procedural de texturas como o processo de criação de texturas por meio de modelos matemáticos e apresentam uma revisão de diferentes métodos para geração de texturas bidimensionais.

Uma formulação simplificada pode ser escrita como:

\[
C=f(P,t),
\]

em que \(P\) representa uma posição espacial, \(t\) representa uma variável temporal quando necessário e \(f\) representa a função procedural responsável pelo resultado.

A abordagem procedural é útil quando se deseja gerar padrões repetitivos, variações ou materiais sem armazenar uma imagem diferente para cada situação. Ela também pode ser combinada com UV coordinates tradicionais.

\section{Vantagens}

Entre as principais vantagens estão a flexibilidade, a possibilidade de animação em tempo real, o reaproveitamento de texturas e a geração de variações sem necessidade de armazenar todos os resultados previamente.

Métodos procedurais também podem reduzir a quantidade de dados necessários para representar determinados tipos de conteúdo. Smelik et al. \cite{smelik2014} destacam que métodos procedurais são atrativos para a geração de grandes quantidades de conteúdo em ambientes virtuais, embora o controle artístico e a interatividade continuem sendo desafios importantes.

\section{Limitações}

A técnica não elimina os problemas tradicionais de mapeamento. Podem ocorrer distorções, descontinuidades, repetição excessiva de padrões e artefatos de filtragem. No caso procedural, funções muito complexas também podem aumentar o custo computacional.

Outro ponto é o controle artístico: uma função procedural pode gerar grande variedade, mas é necessário definir parâmetros que permitam ao artista controlar o resultado desejado.

\section{Custo computacional}

O custo depende principalmente da quantidade de cálculos executados por fragmento, do número de amostragens de textura, dos filtros utilizados e da resolução envolvida. Operações realizadas no fragment shader podem ser executadas muitas vezes, pois um mesmo shader pode processar grande quantidade de fragmentos.

Por isso, uma implementação eficiente deve avaliar quais cálculos realmente precisam ser feitos por fragmento e quais podem ser antecipados ou simplificados.
